\chapter{SOLUSI DAN METODOLOGI}

\section{Solusi yang Ditawarkan}
Solusi yang ditawarkan adalah Welya, aplikasi web dan desktop asisten
akademik berbasis
AI dengan fitur utama:
\begin{enumerate}
  \item Deteksi tugas otomatis dari email kampus (integrasi Gmail);
  \item Pemecahan tugas menjadi subtask beserta estimasi waktu oleh AI;
  \item Tampilan kalender tenggat tugas dan jadwal kuliah (integrasi
    Google Calendar dan Google Tasks);
  \item Pencarian cerdas atas tugas, materi, dan catatan;
  \item Dasbor progres, tenggat, dan prioritas harian.
\end{enumerate}

\section{Target Luaran}
\begin{enumerate}
  \item Aplikasi Welya (web dan desktop) yang berfungsi penuh dan siap
    diuji coba;
  \item Laporan tugas akhir;
  \item Artikel publikasi sesuai ketentuan BAB XIV panduan tugas akhir;
  \item Pendaftaran Hak Kekayaan Intelektual (HKI).
\end{enumerate}

\section{Pihak yang Terlibat}
\begin{enumerate}
  \item Mahasiswa sebagai pengembang produk;
  \item Dosen pembimbing sebidang ilmu dari program studi;
  \item Mahasiswa pengguna sebagai responden uji coba.
\end{enumerate}

\section{Tahapan Pengembangan}
Pengembangan dilakukan secara iteratif dengan mengadaptasi metode
\emph{prototype}, yang memungkinkan perbaikan bertahap berdasarkan umpan
balik pengguna \parencite{aprilisa2024}, melalui tahapan berikut:
\begin{enumerate}
  \item \textbf{Analisis kebutuhan}: identifikasi permasalahan dan
    kebutuhan pengguna melalui studi literatur serta observasi terhadap
    kebiasaan mahasiswa dalam mengelola tugas dan jadwal kuliah;
  \item \textbf{Perancangan sistem}: perancangan arsitektur, basis data,
    alur integrasi layanan Google, dan antarmuka pengguna;
  \item \textbf{Implementasi}: pembangunan aplikasi web dan desktop beserta
    komponen AI untuk analisis email dan penyusunan subtask;
  \item \textbf{Pengujian}: pengujian fungsional serta uji coba kepada
    pengguna terbatas;
  \item \textbf{Evaluasi}: pengukuran kebermanfaatan dan kemudahan
    penggunaan, kemudian perbaikan berdasarkan masukan pengguna.
\end{enumerate}

\section{Deskripsi Rancangan Produk}
Secara garis besar, arsitektur Welya terdiri atas empat komponen utama:
\begin{enumerate}
  \item \textbf{Frontend (aplikasi web dan desktop)}: antarmuka dasbor,
    daftar tugas,
    kalender, dan pencarian yang diakses melalui peramban;
  \item \textbf{Backend}: layanan API yang mengelola autentikasi,
    penyimpanan data tugas dan jadwal, serta orkestrasi proses analisis;
  \item \textbf{Komponen AI}: modul yang memanfaatkan model bahasa besar
    untuk menganalisis isi email, mengekstraksi tugas beserta tenggatnya,
    dan menyusun subtask dengan estimasi waktu pengerjaan;
  \item \textbf{Integrasi layanan Google}: koneksi ke Gmail, Google
    Tasks, dan Google Calendar melalui API resmi dengan otorisasi OAuth
    2.0 \parencite{arianto2025}, sehingga data tetap berada dalam kendali
    pengguna.
\end{enumerate}

Alur kerja sistem dimulai ketika pengguna menghubungkan akun Google-nya.
Sistem membaca email kampus secara berkala, komponen AI menyaring email
yang mengandung penugasan, kemudian tugas yang terdeteksi ditambahkan ke
daftar tugas beserta subtask dan estimasi waktunya. Tenggat tugas dan
jadwal kuliah disinkronkan dengan Google Calendar dan ditampilkan pada
dasbor beserta prioritas harian. Gambaran rancangan yang lebih lengkap
disajikan pada Lampiran 1.

\section{Prosedur Kerja}
\begin{enumerate}
  \item Mengidentifikasi kebutuhan pengguna melalui studi literatur dan
    observasi kebiasaan mahasiswa dalam mengelola tugas;
  \item Merumuskan spesifikasi kebutuhan fungsional dan non-fungsional
    berdasarkan hasil identifikasi kebutuhan;
  \item Merancang arsitektur sistem, skema basis data, dan purwarupa
    antarmuka;
  \item Mengimplementasikan fitur secara bertahap dengan mengacu pada
    praktik pengembangan aplikasi manajemen tugas berbasis web
    \parencite{panjaitan2025}, dimulai dari integrasi akun Google, deteksi
    tugas, hingga dasbor dan pencarian;
  \item Melakukan pengujian fungsional setiap akhir iterasi;
  \item Melaksanakan uji coba kepada responden mahasiswa dan mengukur
    kemudahan penggunaan;
  \item Menganalisis hasil evaluasi dan menyusun laporan.
\end{enumerate}

\section{Rencana Evaluasi}
Evaluasi mencakup tiga aspek. Pertama, pengujian fungsional
(\emph{black-box}) untuk setiap fitur pada akhir tiap iterasi. Kedua,
pengukuran kinerja komponen AI: ketepatan deteksi tugas dari email diukur
dengan \emph{precision} dan \emph{recall} terhadap sampel email kampus
yang dilabeli secara manual, dengan target minimal 80\% untuk keduanya.
Ketiga, kuesioner \emph{System Usability Scale} (SUS) \parencite{dyayu2023}
kepada 20--30 responden mahasiswa; skor SUS di atas 68 menandakan
kebergunaan di atas rata-rata dan menjadi target minimal produk ini.
Hasil evaluasi menjadi dasar perbaikan produk sebelum laporan
akhir disusun.
